\documentclass[12pt]{article}
\usepackage{tabularx}

% Use algorithm and algpseudocode packages so \State and related commands are available

\newcommand{\duedate}{---}
\newcommand{\assignment}{MCTS (Lec5) Summary} % Change to "Problem Set X"

% Change the following to your name and UNI.
\newcommand{\name}{Aksel Kretsinger-Walters, adk2164}
\newcommand{\email}{adk2164@columbia.edu}

% NOTE: Defining collaborators is optional; to not list collaborators, comment out the
% line below. Maximum of two collaborators per problem set.
%\newcommand{\collaborators}{Vig Vigerton (\texttt{UNI}), Alice Bob (\texttt{UNI})}
% No collaborators on PS 0

\makeatletter
\def\input@path{{../}{../../}{../../../}} % add as many parents as you need
\makeatother
\input{pset_template_RL.tex} %% DO NOT CHANGE THIS LINE

\begin{document}
\psetheader %% DO NOT CHANGE THIS LINE

\section*{Lecture Summary: Monte Carlo Tree Search (MCTS)}

\subsection*{Overview}
Monte Carlo Tree Search (MCTS) is a simulation–based planning method that combines
\emph{tree search} with \emph{multi-armed bandit} ideas.
It is designed for settings where:
\begin{itemize}
		\item a simulator exists and can reset to arbitrary states,
		\item rewards are sparse (often only at terminal states),
		\item the branching factor is too large for minimax or $\alpha$–$\beta$ pruning,
		\item exhaustive search is infeasible.
\end{itemize}
MCTS adaptively focuses computation on the most promising trajectories and forms the
core of modern game-playing systems such as MoGo, AlphaGo, AlphaZero, and MuZero.

\subsection*{Historical Motivation}
Early successes in Computer Go revealed that:
\begin{itemize}
		\item high branching factors make classical game-tree search intractable,
		\item playout (simulation)–based evaluation is effective when a simulator is available,
		\item bandit algorithms can guide search toward promising moves.
\end{itemize}
The key idea: \emph{Tree Search + Bandits = MCTS}.
Each node in the tree is treated like a bandit problem, and previously successful
actions are explored more deeply.

\subsection*{Upper Confidence Bounds in Trees (UCT)}
The most widely used MCTS variant is UCT (Kocsis \& Szepesvári, 2006).

\paragraph{Goal.}
Find a path from the root to a leaf with highest expected reward in a tree of depth $D$
and branching factor $K$.

\paragraph{UCT Score.}
At a node $j$, for each child $k = 1,\dots,K$, UCT selects the child maximizing
\[
		B_t(k) \;=\; \hat{\mu}_{k,t}
		+ c \sqrt{\frac{\log(n_{j,t})}{n_{k,t}}}.
\]
Here:
\begin{itemize}
		\item $\hat{\mu}_{k,t}$ is the empirical mean reward of subtree $k$,
		\item $n_{k,t}$ is the number of times child $k$ has been visited,
		\item $n_{j,t}$ counts visits to node $j$,
		\item $c > 0$ controls exploration.
\end{itemize}
This mirrors the classical UCB rule but is applied hierarchically across the tree.

\subsection*{Phases of the MCTS Algorithm}
One iteration of MCTS consists of:

\paragraph{1. Selection.}
Starting at the root, recursively choose the child with highest UCT score until reaching a leaf.

\paragraph{2. Expansion.}
If the leaf corresponds to a previously visited state and actions remain unexplored,
create new child nodes (often using a simulator or \texttt{deepcopy}).

\paragraph{3. Simulation (Rollout).}
Perform a random or heuristic playout from the expanded node until termination,
generating a Monte-Carlo estimate of the value.

\paragraph{4. Backpropagation.}
Propagate the final reward back up the path, updating visit counts and empirical means.

\subsection*{Theoretical Guarantees}
\paragraph{Consistency.}
As the number of simulations grows:
\[
		\hat{\mu}_{i,t} \to \mu_i \quad \text{for each leaf } i,
\]
and UCT eventually selects the optimal path with probability $1$.

\paragraph{Finite-Time Performance.}
Analysis is subtle because:
\begin{itemize}
		\item leaf nodes behave like stationary bandits,
		\item internal nodes behave like \emph{non-stationary} bandits (their estimates drift
				as deeper nodes become better understood),
		\item UCB analysis can be adapted to mild drift conditions.
\end{itemize}
A regret bound of order $O(K \log T)$ can be obtained, but additive constants may be
\emph{doubly exponential} in tree depth $D$.
Thus finite-time performance may be extremely poor in worst-case trees.

\subsection*{Practical Behavior}
In some adversarial tree structures, UCT may spend exponentially many simulations
on suboptimal branches before discovering the optimal one.
This is because UCB pulls suboptimal arms only $O(\log n)$ times, delaying discovery
when high-reward leaves are deep and rare.

\subsection*{Modern Extensions}
\paragraph{AlphaZero.}
Combines UCT with:
\begin{itemize}
		\item a value network $V(s)$,
		\item a policy network $\pi(s,a)$,
		\item mixed scores that combine UCT bonus, neural value predictions, and policy priors,
		\item bootstrapping or full Monte-Carlo rollouts.
\end{itemize}

\paragraph{MuZero.}
Learns its own dynamics model to simulate rollouts in a latent space and uses value
bounds to stabilize planning.

\paragraph{General Theme.}
MCTS scales to large state spaces by coupling:
\begin{itemize}
		\item bandit-style exploration,
		\item function approximation (values, policies, or dynamics),
		\item selective lookahead search.
\end{itemize}
It remains the foundation for state-of-the-art reasoning and planning systems.
\end{document}
